\paragraph{Remark.}
The separation of centroid-in and centroid-out nodes is essential for the
validity of the forward dynamic-programming loading algorithm. Without
this separation, egress links could break the acyclic structure and violate the topological order implied by time.

\begin{definition}[Links]
The link set $\mathcal{A}$ contains the following link types:
\begin{enumerate}
\item \textbf{Access links (time-windowed):}
      for an OD pair $q=(s_o,s_d)$ and departure interval $t$, access links connect the
      origin centroid-in node $(s_o,t)^{\mathrm{in}}$ to \emph{selected}
      \emph{departure} event nodes at the origin stop,
      \[
      (s_o,t)^{\mathrm{in}} \rightarrow (s_o,\tau)^{\mathrm{dep}},
      \]
only for departure times $\tau$ satisfying
\[
\tau \in \bigl[a_t-\Delta_{\max}^{\mathrm{access}},\; b_t+\Delta_{\max}^{\mathrm{access}}\bigr].
\]

\item \textbf{Ride links:}
      for each trip, ride links connect a departure event at the current stop to an
      arrival event at the next stop,
      \[
      (s_i,\tau_i^{\mathrm{dep}})^{\mathrm{dep}} \rightarrow (s_{i+1},\tau_{i+1}^{\mathrm{arr}})^{\mathrm{arr}}.
      \]

\item \textbf{Dwell/continue links (within-trip at a stop):}
      for each trip and each stop event, a dwell link connects the arrival and departure
      nodes at the same stop,
      \[
      (s,\tau^{\mathrm{arr}})^{\mathrm{arr}} \rightarrow (s,\tau^{\mathrm{dep}})^{\mathrm{dep}}.
      \]
      In the implementation, strictly positive dwell time is enforced by regularizing
      timetable records where arrival equals departure (using a small minimum dwell time).

\item \textbf{Transfer links (inter-line, time-windowed):}
      transfer links connect an \emph{arrival} event node to a later \emph{departure} event node
      at the same stop and represent waiting and switching \emph{between different lines}. A transfer link
      \[
      (s,\tau_1)^{\mathrm{arr}} \rightarrow (s,\tau_2)^{\mathrm{dep}},
      \qquad\text{with}\qquad \tau_2>\tau_1
      \]
      is included only if
      \begin{enumerate}
      \item the two event nodes belong to \emph{different lines}, and
      \item the waiting time satisfies $\tau_2-\tau_1 \le \Delta_{\max}^{\mathrm{transfer}}$.
      \end{enumerate}
      No transfer links are created between events of the same line. In particular,
      transfer links are never created for $\tau_2=\tau_1$.

\item \textbf{Egress links (global graph, destination-gated by masking):}
      the global graph contains egress links from every arrival event node at stop $s$
      to the centroid-out node $s^{\mathrm{out}}$:
      \[
      (s,\tau)^{\mathrm{arr}} \rightarrow s^{\mathrm{out}}.
      \]
      For a given OD pair $q=(s_o,s_d)$, only the destination centroid-out node $s_d^{\mathrm{out}}$
      is enabled as a terminal node; all other egress options are disabled by a destination-dependent mask.
\end{enumerate}
\end{definition}

Each link $\ell\in\mathcal{A}$ is associated with:
\begin{itemize}
\item a generalized cost $c_\ell$,
\item a capacity $u_\ell$ for ride links (stored in the data structure, not yet used in costs).
\end{itemize}

\begin{assumption}[Acyclic time-expanded network]
All links in the time-expanded network are consistent with a strict
topological order. All \emph{event-to-event} links strictly increase
clock time. In addition, centroid-in nodes $(s,t)^{\mathrm{in}}$ are
placed \emph{before} all event nodes (conceptually at $-\infty$), and
centroid-out nodes $s^{\mathrm{out}}$ are placed \emph{after} all event
nodes (conceptually at $+\infty$). With this convention, access and
egress links always respect the global ordering.

With the event split, feasible movements follow the pattern
\[
(s,t)^{\mathrm{in}}
\rightarrow (s,\tau)^{\mathrm{dep}}
\rightarrow (s',\tau')^{\mathrm{arr}}
\rightarrow (s',\tau'')^{\mathrm{dep}}
\rightarrow \cdots
\rightarrow (s_d,\bar{\tau})^{\mathrm{arr}}
\rightarrow s_d^{\mathrm{out}}.
\]
Transfer links satisfy $\tau_2>\tau_1$ and dwell links are enforced to be strictly forward in time
through a minimum dwell regularization. Therefore the directed graph is acyclic and admits
a topological ordering, which allows value functions and flows to be computed using dynamic programming.
\end{assumption}

